\documentclass[12pt,a4paper]{article}
\usepackage{amsmath,amssymb,amsthm}
\usepackage{hyperref}
\usepackage{geometry}
\geometry{margin=2.5cm}
\usepackage{enumitem}

\newtheorem{theorem}{Theorem}[section]
\newtheorem{lemma}[theorem]{Lemma}
\newtheorem{corollary}[theorem]{Corollary}
\newtheorem{proposition}[theorem]{Proposition}
\theoremstyle{definition}
\newtheorem{definition}[theorem]{Definition}
\theoremstyle{remark}
\newtheorem{remark}[theorem]{Remark}

\title{ANITA Upgoing Events from\\
Substrate-Mediated Particle Production}

\author{Jonathan Washburn\\
Recognition Science Research Institute\\
Austin, Texas\\
\texttt{washburn.jonathan@gmail.com}}

\date{March 2026}

\begin{document}
\maketitle

\begin{abstract}
We analyze the ANITA anomalous upgoing events within Recognition
Science (RS).  The ANITA balloon experiment detected two
cosmic-ray-like events at energies $\sim 0.6~\mathrm{EeV}$, apparently
emerging from below the Antarctic ice at steep zenith angles
($\sim 25^\circ$--$35^\circ$ below the horizon)~\cite{ANITA2016,
ANITA2018}.  Standard Model neutrinos at these energies have
interaction lengths much shorter than the Earth's diameter, making
upgoing $\nu$-induced air showers effectively impossible.  Proposed BSM
explanations include stau sleptons, sterile neutrinos, and sphaleron
transitions.  RS offers a qualitative interpretation: the recognition ledger can
transiently convert stored coherence energy into particle showers at
high-density nodes of the $\varphi$-lattice, producing upgoing
air-shower-like signatures without requiring upgoing neutrinos or
BSM particles.  This mechanism is a hypothesis; no rate prediction
is made here.  RS predicts that no BSM particle (stau slepton, sterile
neutrino) will account for the events, and that any confirmed upgoing
events will be rare and clustered near $\varphi$-rung energies.  The
PUEO experiment (2022 flight) and ARA will provide further data.
\end{abstract}

%% ============================================================
\section{Recognition Science Framework}
\label{sec:framework}
%% ============================================================

The ANITA analysis relies on the RS substrate structure, which is
uniquely fixed by the following cost axioms.

\begin{definition}[RS Cost Functional]
For $x > 0$: $J(x) = \tfrac{1}{2}(x + x^{-1}) - 1$.
\end{definition}

\noindent\textbf{Axiom A1 (Normalization):} $J(1) = 0$.\\
\textbf{Axiom A2 (Recognition Composition Law, RCL):}
$J(xy)+J(x/y)=2J(x)J(y)+2J(x)+2J(y)$ for all $x,y>0$.\\
\textbf{Axiom A3 (Calibration):} $J''_{\log}(0) = 1$, where
$J_{\log}(t) := J(e^t)$.

\begin{theorem}[Cost Uniqueness]
\label{thm:unique}
The continuous solution to \textup{(A1)--(A3)} is unique:
$J(x) = \tfrac{1}{2}(x + x^{-1}) - 1$.
\end{theorem}

\begin{proof}[Proof sketch]
Set $H(t) = J_{\log}(t)+1$.  Axiom A2 transforms into the d'Alembert
equation $H(s+t)+H(s-t)=2H(s)H(t)$.  By the Acz\'el classification
theorem~\cite{Aczel1966}, the unique continuous solution with $H(0)=1$
and $H''(0)=1$ is $H(t)=\cosh t$, giving
$J(x)=\tfrac{1}{2}(x+x^{-1})-1$.
\end{proof}

\noindent The recognition ledger stores coherence energy at nodes of the
discrete $\varphi$-lattice (with $\varphi=(1+\sqrt{5})/2$ forced by RCL
self-similarity).  Space is a $\mathbb{Z}^3$ lattice ($D=3$ forced by
$2^D=8$); the substrate couples to matter density.  At high-density
nodes, transient conversion of stored coherence energy into particle
showers is the proposed mechanism for the ANITA events.

%% ============================================================
\section{Introduction}
\label{sec:intro}
%% ============================================================

Recognition Science~\cite{RS_Axioms} (RS) provides a substrate
coherence framework for interpreting rare high-energy events without
requiring beyond-Standard-Model particles.

The Antarctic Impulsive Transient Antenna (ANITA) is a balloon-borne
radio detector designed to observe ultra-high-energy cosmic neutrinos
via Askaryan radiation in Antarctic ice.  During ANITA-I (2006) and
ANITA-III (2014), two anomalous events were detected:

\begin{itemize}[nosep]
  \item Event AAE061228 (ANITA-I): $E \approx 0.6$~EeV, zenith angle
  $\sim 117^\circ$ (upgoing at $\sim 27^\circ$ below
  horizon)~\cite{ANITA2016}.
  \item Event AAE141220 (ANITA-III): $E \approx 0.56$~EeV, zenith
  angle $\sim 125^\circ$ ($\sim 35^\circ$ below
  horizon)~\cite{ANITA2018}.
\end{itemize}

At $E \sim 0.6$~EeV, the $\nu N$ cross section is
$\sigma_{\nu N} \sim 10^{-31}~\mathrm{cm}^2$, giving an interaction
length $\lambda \sim 10~\mathrm{km}$ in rock.  The Earth's chord
lengths at the observed angles are $\sim 5000$--$7000~\mathrm{km}$,
yielding survival probability $e^{-L/\lambda} \sim 10^{-200}$.
Standard neutrinos cannot produce these events.

%% ============================================================
\section{RS Analysis}
\label{sec:analysis}
%% ============================================================

\begin{theorem}[Standard Neutrinos Cannot Produce These Events]
\label{thm:nu-impossible}
At energy $E \sim 0.6$~EeV, the neutrino interaction length in
rock is $\lambda_\nu \sim 10$~km, while the Earth chord at the
observed zenith angles ($\sim 117$--$125^\circ$) is
$L \sim 5000$--$7000$~km.  The neutrino survival probability is:
\begin{equation}
  P_{\mathrm{survive}} = e^{-L/\lambda_\nu} \lesssim e^{-500}
  \approx 10^{-217}.
\end{equation}
This is effectively zero; standard neutrinos cannot produce
Earth-traversing showers at these energies.
\end{theorem}

\begin{proposition}[RS Substrate Interpretation]
\label{prop:substrate}
The recognition ledger stores coherence energy at nodes of the
$\varphi$-lattice.  At high-density nodes (where matter density
enhances substrate coherence near the Antarctic crust), the
substrate can transiently convert stored coherence energy into
particle showers at energies $E \sim E_{\mathrm{coh}} \cdot \varphi^N$
for appropriate rung $N$.  The resulting shower propagates upward,
producing a radio signal indistinguishable from a downgoing air shower
but with upgoing geometry.  This is a qualitative hypothesis; no
quantitative rate is derived in the present work.
\end{proposition}

\begin{remark}
No stau sleptons, sterile neutrinos, or other BSM particles are
required under the RS substrate interpretation.  The key prediction
distinguishing RS from BSM explanations is the \emph{absence} of
a correlated flux of upgoing muons or tau-induced showers at IceCube
from the same directions.
\end{remark}

%% ============================================================
\section{Predictions}
\label{sec:predict}
%% ============================================================

\begin{enumerate}
  \item \textbf{PUEO rate.}  The Payload for Ultrahigh Energy
  Observations (PUEO), ANITA's successor, should observe upgoing
  events at a rate consistent with rare substrate production events
  ($\lesssim 1$ per flight), not a continuous neutrino flux.

  \item \textbf{No correlated flux.}  No corresponding flux of upgoing
  muons or taus should be observed by IceCube or ARA from the same
  directions, since the events are not produced by neutrinos.

  \item \textbf{Energy clustering.}  Substrate production events
  should cluster at energies corresponding to $\varphi$-rungs, not at
  arbitrary energies.

  \item \textbf{Geographic correlation.}  Events should preferentially
  occur at locations with enhanced substrate coherence (geological
  density anomalies, tectonic plate boundaries).
\end{enumerate}

%% ============================================================
\section{Conclusion}
\label{sec:conclusion}
%% ============================================================

The ANITA upgoing events---two events with energies $\sim 0.6$~EeV
at steep zenith angles---cannot be produced by standard upgoing
neutrinos (survival probability $\lesssim 10^{-200}$).  The RS
substrate interpretation offers a qualitative mechanism via coherence
energy release at high-density $\varphi$-lattice nodes, without
requiring BSM particles.  A quantitative derivation of the event
rate within RS remains as future work.  The PUEO and ARA experiments
will determine whether the ANITA events recur at a predictable rate.

%% ============================================================
\begin{thebibliography}{99}

\bibitem{ANITA2016}
P.~W.~Gorham et al., Phys.\ Rev.\ Lett.\ \textbf{117}, 071101 (2016).

\bibitem{ANITA2018}
P.~W.~Gorham et al., Phys.\ Rev.\ Lett.\ \textbf{121}, 161102 (2018).

\bibitem{RS_Axioms}
J.~Washburn, ``The Algebra of Reality: A Recognition Science Derivation
of Physical Law,'' \emph{Axioms} \textbf{15}(2), 90 (2026).
\texttt{doi:10.3390/axioms15020090}

\bibitem{Aczel1966}
J.~Acz\'el, \emph{Lectures on Functional Equations and Their
Applications} (Academic Press, 1966).

\end{thebibliography}

\end{document}
